%%%%%%%%%%%%%%%%%%%%%%%%%%%%%%%%%%%%%%%%%%%%%%%%%%%%%%%%%%%%%%%%%%%%%%%%%%%%%%%%%%%%
% Math notation 
%%%%%%%%%%%%%%%%%%%%%%%%%%%%%%%%%%%%%%%%%%%%%%%%%%%%%%%%%%%%%%%%%%%%%%%%%%%%%%%%%%%%

%%% Graphs 
% \newcommand{\CI}{\mathrel{\perp\mspace{-10mu}\perp}}
% \newcommand{\olineo}{\circ\mspace{-7mu}-\mspace{-7mu}\circ}
% \newcommand{\orightarrow}{\circ\!\!\rightarrow}
% \newcommand{\oleftarrow}{\leftarrow\!\!\circ}
% \newcommand{\sleftarrow}{\leftarrow\mspace{-7.5mu}*}
% \newcommand{\srightarrow}{*\mspace{-7.5mu}\rightarrow}
% \newcommand{\slineo}{*\mspace{-6mu}-\mspace{-6mu}\circ}
% \newcommand{\olines}{\circ\mspace{-6mu}-\mspace{-6mu}*}
% \newcommand{\slines}{*\mspace{-6mu}-\mspace{-6mu}*}
% \newcommand{\directedpath}{--\mspace{-5mu}\rightarrow}
% \newcommand{\divergentpath}{\leftarrow \mspace{-5mu} * \mspace{-5mu}\rightarrow}    
% \newcommand{\G}{\mathcal{G}}
% \newcommand{\Gin}[1]{\mathcal{G}_{\overline{#1}}}
% \newcommand{\Gout}[1]{\mathcal{G}_{\underline{#1}}}
% \newcommand{\Gio}[2]{\mathcal{G}_{\overline{#1}\underline{#2}}}


% %%% Notations 
% \newcommand{\omf}[2]{\overline{#1}^{#2}}
% \newcommand{\mf}[1]{\mathbf{#1}}
% \newcommand{\pre}[2]{\text{pre}_{#1}(#2)}
% \newcommand{\prev}[1]{\text{pre}(#1)}
% \newcommand{\ogeq}[2]{\mf{#1}^{\geq #2}}
% \newcommand{\oleq}[1]{\mf{#1}^{i \leq}}
% \newcommand{\lr}[1]{\left(#1\right)}
% \newcommand{\lrp}[1]{\left\{#1\right\}}
% \newcommand{\lrb}[1]{\left[#1\right]}
% \newcommand{\doo}[1]{\operatorname{do}(#1)}

% %%% Math  
% \newcommand{\domain}[1]{\mathfrak{D}_{#1}}
% \newcommand{\support}[1]{\mathfrak{S}_{#1}}
% \newcommand{\cupdot}{\mathbin{\mathaccent\cdot\cup}}
% \newcommand{\inner}[2]{\left\langle #1, #2 \right\rangle}
% \newcommand{\notimplies}{\mathrel{{\ooalign{\hidewidth$\not\phantom{=}$\hidewidth\cr$\implies$}}}}
% \newcommand{\abs}[1]{\left\vert #1 \right\vert}
% \newcommand{\norm}[1]{\left\lVert#1\right\rVert}
% \newcommand{\I}[2]{\mathbb{I}_{#1}(#2)}
% \newcommand{\deriv}[1]{\frac{\partial}{\partial #1}}
% \newcommand{\deriveval}[2]{(\partial / \partial #1)\vert_{#2}}
% \newcommand{\integral}[3]{\int_{#1} #2 \;d[#3]}

% %%% Estimator  
% \newcommand{\mean}[1]{\mathbb{E}[#1]}
% \newcommand{\meani}[2]{\mathbb{E}_{#1}[#2]}
% \newcommand{\variance}[1]{\mathbb{V}[#1]}
% \newcommand{\variancei}[2]{\mathbb{V}_{#1}[#2]}
% \newcommand{\bigmean}[1]{\mathbb{E}\left[#1\right]}
% \newcommand{\bigmeani}[2]{\mathbb{E}_{#1}\left[#2\right]}
% \newcommand{\bigvariance}[1]{\mathbb{V}\left[#1\right]}
% \newcommand{\bigvariancei}[2]{\mathbb{V}_{#1}\left[#2\right]}
% \newcommand{\ith}[2]{#1_{(#2)}}



%%%%%%%%%%%%%%%%%%%%%%%%%%%%%%%%%%%%%%%%%%%%%%%%%%%%%%%%%%%%%%%%%%%%%%%%%%%%%%%%%%%%
% Setting 
%%%%%%%%%%%%%%%%%%%%%%%%%%%%%%%%%%%%%%%%%%%%%%%%%%%%%%%%%%%%%%%%%%%%%%%%%%%%%%%%%%%%
\let\oldnl\nl% Store \nl in \oldnl
\newcommand{\nonl}{\renewcommand{\nl}{\let\nl\oldnl}}% Remove line number for one line
\definecolor{opaquegray}{RGB}{192, 192, 192}
\definecolor{afblue}{RGB}{0,0,102}
\definecolor{mypink}{RGB}{255,51,153}

% \newtheorem{definition}{Definition}
% \newtheorem{theorem}{Theorem}
\newtheorem{assumption}{Assumption}
% \newtheorem{example}{Example}
% \newtheorem{lemma}{Lemma}
% \newtheorem{conjecture}{Conjecture}
% \newtheorem{proposition}{Proposition}
% \newtheorem{corollary}{Corollary}
% \newtheorem{preliminary}{Preliminary}
% \newtheorem{remark}{Remark}
% \newtheorem{axiom}{Axiom}
% \newtheorem{principle}{Principle}
% \newtheorem{task}{Task}

% Example for Definition
\tcolorboxenvironment{definition}{
  colback=gray!5!white,
  colframe=black,
  % colback=blue!5!white,
  % colframe=blue!80!black,
  fonttitle=\bfseries,
  boxrule=0.5pt, % Adjust this value for thinner or thicker boundary
  breakable,
  % sharp corners,
}

\tcolorboxenvironment{theorem}{
  colback=gray!5!white,
  colframe=black,
  % colback=blue!5!white,
  % colframe=blue!80!black,
  fonttitle=\bfseries,
  boxrule=0.5pt, % Adjust this value for thinner or thicker boundary
  breakable,
  % sharp corners,
}

\tcolorboxenvironment{corollary}{
  colback=gray!5!white,
  colframe=black,
  % colback=blue!5!white,
  % colframe=blue!80!black,
  fonttitle=\bfseries,
  boxrule=0.5pt, % Adjust this value for thinner or thicker boundary
  breakable,
  % sharp corners,
}

\tcolorboxenvironment{lemma}{
  colback=gray!5!white,
  colframe=black,
  % colback=blue!5!white,
  % colframe=blue!80!black,
  fonttitle=\bfseries,
  boxrule=0.5pt, % Adjust this value for thinner or thicker boundary
  breakable,
  % sharp corners,
}

\tcolorboxenvironment{proposition}{
  colback=gray!5!white,
  colframe=black,
  % colback=blue!5!white,
  % colframe=blue!80!black,
  fonttitle=\bfseries,
  boxrule=0.5pt, % Adjust this value for thinner or thicker boundary
  breakable,
  % sharp corners,
}

\tcolorboxenvironment{proof}{
  colback=white,
  colframe=black,
  % colback=blue!5!white,
  % colframe=blue!80!black,
  fonttitle=\bfseries,
  boxrule=0.5pt, % Adjust this value for thinner or thicker boundary
  breakable,
  % sharp corners,
}

\tcolorboxenvironment{example}{
  colback=white,
  colframe=black,
  % colback=blue!5!white,
  % colframe=blue!80!black,
  fonttitle=\bfseries,
  boxrule=0.5pt, % Adjust this value for thinner or thicker boundary
  breakable,
  % sharp corners,
}

\tcolorboxenvironment{remark}{
  colback=gray!5!white,
  colframe=black,
  % colback=blue!5!white,
  % colframe=blue!80!black,
  fonttitle=\bfseries,
  boxrule=0.5pt, % Adjust this value for thinner or thicker boundary
  breakable,
  % sharp corners,
}

\tcolorboxenvironment{assumption}{
  colback=white,
  colframe=black,
  % colback=blue!5!white,
  % colframe=blue!80!black,
  fonttitle=\bfseries,
  boxrule=0.5pt, % Adjust this value for thinner or thicker boundary
  breakable,
  % sharp corners,
}

% % Example for Theorem
% \tcolorboxenvironment{theorem}{
%   colback=red!5!white,
%   colframe=red!80!black,
%   fonttitle=\bfseries,
%   boxrule=0.5pt, % Adjust this value for thinner or thicker boundary
%   % sharp corners,
%   breakable,
% }

% % Example for Lemma
% \tcolorboxenvironment{lemma}{
%   colback=orange!5!white,
%   colframe=orange!80!black,
%   fonttitle=\bfseries,
%   boxrule=0.5pt,
%   % sharp corners,
%   breakable,
% }

% % Example for Example (demonstration)
% \tcolorboxenvironment{example}{
%   colback=green!5!white,
%   colframe=green!80!black,
%   fonttitle=\bfseries,
%   boxrule=0.5pt,
%   % sharp corners,
%   breakable,
% }
